\documentclass[11pt,a4paper]{article}
\usepackage[margin=24mm]{geometry}
\usepackage{amsmath,amssymb,booktabs,array,hyperref,listings,microtype}
\usepackage[T1]{fontenc}
\usepackage{lmodern}
\hypersetup{colorlinks=true,linkcolor=blue,urlcolor=blue,pdftitle={Complete MTK Hydropower SMIB Model}}
\lstdefinelanguage{Julia}{keywords={using,begin,end,function,return,if,else,elseif,for,while,module,export,include,struct},sensitive=true,comment=[l]\#,morestring=[b]"}
\lstset{language=Julia,basicstyle=\ttfamily\small,breaklines=true,frame=single,columns=fullflexible,showstringspaces=false}

\title{Complete MTK Hydropower SMIB Model\\\large Refractoring based on state-of-the art work flow}
\author{OpenHPLjl development note}
\date{20 September 2026}

\begin{document}
\maketitle

\begin{abstract}
This document records the package commit that turns the OpenHPLjl component catalogue into a complete reduced hydropower single-machine infinite-bus (SMIB) model in ModelingToolkit.jl. The model couples hydraulic storage and water-column dynamics, turbine power conversion, shaft inertia, synchronous-generator angle dynamics, an infinite grid, frequency measurement, and a droop governor. The implementation follows a layered workflow: physical law, connector semantics, reusable component, hierarchical composition, structural compilation, consistent initialization, numerical simulation, and verification.
\end{abstract}

\section{Commit objective}
The objective of this commit is to move from individually tested components to one connected nonlinear plant model that can be compiled and simulated as a single ModelingToolkit system. The model is deliberately reduced, but every major energy domain required for governor and frequency-response studies is represented explicitly.

The top-level model is exported as \texttt{ReducedSMIB}. It is intended to be the baseline system for subsequent guide-vane disturbance tests, governor tuning, AGC extensions, FCR studies, operating-point linearization, Bode/Nyquist analysis, and plant-specific parameterization.

\section{System architecture}
The physical and control architecture is

\begin{verbatim}
Upstream reservoir
      |
      v
   Headrace
      |
      v
  Surge tank
      |
      v
   Penstock
      |
      v
Shaft-coupled turbine ---> Shaft ---> SMIB generator ---> Infinite grid
          ^                   |
          |                   v
          +--- Droop governor <- Frequency sensor
\end{verbatim}

The system contains eleven named subsystems:

\begin{center}
\begin{tabular}{ll}
\toprule
Subsystem & Role \\
\midrule
\texttt{upstream} & constant-head upstream boundary \\
\texttt{headrace} & hydraulic inertia / momentum balance \\
\texttt{surge} & local hydraulic storage \\
\texttt{penstock} & downstream water-column inertia \\
\texttt{turbine} & hydraulic-to-mechanical conversion \\
\texttt{tailwater} & constant-head downstream boundary \\
\texttt{shaft} & rotating inertia and damping \\
\texttt{generator} & SMIB electromechanical conversion \\
\texttt{grid} & infinite-frequency electrical boundary \\
\texttt{frequency\_sensor} & non-loading speed measurement \\
\texttt{governor} & first-order droop control \\
\bottomrule
\end{tabular}
\end{center}

\section{Connector organization}
The implementation separates physical connectors from causal control signals.

\subsection{Hydraulic connector}
The hydraulic interface uses head as the potential variable and volumetric flow as the conserved flow variable:
\[
H \qquad Q.
\]
The sign convention is flow positive into a component. Connected hydraulic flow variables therefore sum to zero.

\subsection{Rotational connector}
The mechanical interface uses
\[
\omega \qquad \tau,
\]
where angular speed is the potential variable and torque is the conserved flow variable.

\subsection{Electrical power connector}
The reduced grid interface uses
\[
\omega_e \qquad P,
\]
where electrical angular frequency is the across variable and active power is the conserved flow variable.

\subsection{Signal connectors}
Control and measurement signals use causal connectors rather than physical flow/potential semantics:
\begin{verbatim}
SignalPlug  ->  SignalSocket
   output          input
\end{verbatim}
This allows the turbine gate command to be driven by a governor, AGC block, test signal, or other controller without changing the turbine physics.

\section{Hydraulic model}
\subsection{Headwater and tailwater}
The two reservoir boundaries prescribe hydraulic head,
\[
H_{up}=120\;\mathrm{m}, \qquad H_{down}=20\;\mathrm{m}
\]
for the default case, while the connected network determines boundary flows.

\subsection{Rigid water columns}
Both headrace and penstock use
\[
\dot Q =
\frac{gA}{L}
\left(
H_{in}-H_{out}-h_f(Q)
\right).
\]
The complete-model baseline uses \texttt{:none} friction so that the default operating point can be constructed exactly before additional hydraulic-loss complexity is introduced.

\subsection{Surge tank}
The surge tank follows
\[
A_s\dot H_s = Q_{in}+Q_{out}.
\]
At the nominal condition, equal steady flow passes through the headrace and penstock, so \(\dot H_s=0\).

\section{Turbine and shaft coupling}
The new \texttt{ShaftCoupledTurbine} closes the hydraulic-mechanical link.

Turbine head is
\[
H_t=H_{in}-H_{out},
\]
guide-vane opening is supplied by a causal input,
\[
y=u_{gate},
\]
and the baseline flow law is
\[
Q=K_q y\sqrt{H_t}.
\]

Mechanical power is
\[
P_h=\rho g\eta_t QH_t.
\]
With the package convention that torque is positive into a component, the generating turbine exports mechanical power through
\[
\tau_t\omega=-P_h.
\]

The shaft supplies the rotational inertia:
\[
J\dot\omega =
\tau_{drive}
+
\tau_{load}
-
D(\omega-\omega_0).
\]

\section{Reduced SMIB generator}
The new \texttt{SMIBGenerator} adds a rotor-angle state relative to the infinite bus:
\[
\dot\delta=\omega_m-\omega_g.
\]

Mechanical and electrical power satisfy
\[
P_m=\tau_g\omega_m,
\]
\[
P_e=\eta_gP_m,
\]
and the reduced infinite-bus transfer relation is
\[
P_e=P_{max}\sin\delta.
\]

Generated electrical power leaves the generator component:
\[
P_{grid}=-P_e.
\]

Rotational inertia remains in \texttt{LumpedShaft}; it is not duplicated in the generator.

\section{Measurement and governor loop}
The \texttt{FrequencySensor} reads shaft speed without mechanically loading the shaft:
\[
\tau_{sensor}=0,
\]
and reports
\[
f=\frac{\omega}{2\pi}.
\]

The droop governor computes
\[
y_{cmd}
=
y_0+\frac{f_{ref}-f}{R},
\]
and applies first-order actuator/governor dynamics:
\[
T_g\dot y=y_{cmd}-y.
\]

The complete feedback path is therefore
\[
\omega
\rightarrow
f
\rightarrow
\text{droop governor}
\rightarrow
y
\rightarrow
Q
\rightarrow
P_h
\rightarrow
\omega.
\]

\section{Consistent default operating point}
The default values used by \texttt{ReducedSMIB} are:

\begin{center}
\begin{tabular}{lll}
\toprule
Quantity & Symbol & Default \\
\midrule
Upstream head & \(H_{up}\) & 120 m \\
Tailwater head & \(H_{down}\) & 20 m \\
Initial flow & \(Q_0\) & 10 m$^3$/s \\
Gate opening & \(y_0\) & 1.0 \\
Grid frequency & \(f_0\) & 50 Hz \\
Turbine efficiency & \(\eta_t\) & 0.90 \\
Generator efficiency & \(\eta_g\) & 0.98 \\
Generator transfer limit & \(P_{max}\) & 10 MW \\
Surge area & \(A_s\) & 500 m$^2$ \\
Governor time constant & \(T_g\) & 0.2 s \\
Droop coefficient & \(R\) & 2.5 Hz/pu gate \\
\bottomrule
\end{tabular}
\end{center}

The initial turbine head is
\[
H_{t0}=H_{up}-H_{down}=100\;\mathrm{m}.
\]

The turbine coefficient is constructed from the requested nominal flow,
\[
K_{q0}
=
\frac{Q_0}
{y_0\sqrt{H_{t0}}}.
\]

Initial turbine power is
\[
P_{h0}
=
\rho g\eta_tQ_0H_{t0}
=
8.829\;\mathrm{MW}.
\]

Initial electrical power is
\[
P_{e0}
=
\eta_gP_{h0}
=
8.65242\;\mathrm{MW}.
\]

The initial rotor angle is computed from
\[
\delta_0
=
\sin^{-1}
\left(
\frac{P_{e0}}{P_{max}}
\right),
\]
instead of being selected as an arbitrary guess.

At the default condition the governor sees 50 Hz and returns \(y=y_0\), water-column flows are equal, the surge level is stationary, and mechanical/electrical power are balanced.

\section{Top-level ModelingToolkit composition}
The complete system is assembled from components and \texttt{connect} statements. The top-level system does not rewrite the internal equations of individual components.

\begin{lstlisting}
@named plant = ReducedSMIB()

compiled = mtkcompile(plant)

prob = ODEProblem(
    compiled,
    [],
    (0.0, 5.0),
)

sol = solve(
    prob,
    Rodas5P();
    abstol = 1e-8,
    reltol = 1e-8,
)
\end{lstlisting}

The workflow deliberately preserves two model views:
\begin{itemize}
\item the hierarchical uncompiled system for component, port, signal, and future analysis-point inspection;
\item the compiled system for structural simplification, initialization, and numerical solution.
\end{itemize}

\section{Files introduced by the commit}
The principal new implementation files are:

\begin{verbatim}
src/Measurements/FrequencySensor.jl
src/Controls/DroopGovernor.jl
src/Turbines/ShaftCoupledTurbine.jl
src/Electrical/Generators/SMIBGenerator.jl
src/Systems/ReducedSMIB.jl
examples/reduced_smib.jl
\end{verbatim}

The package entry point is updated to include and export these components, and \texttt{test/runtests.jl} is extended with a complete plant compile-and-simulate test.

\section{Verification strategy}
The commit follows a verification ladder rather than only checking whether Julia code parses.

\subsection*{Component verification}
Existing tests continue to check reservoir mass balance, rigid-pipe momentum response, surge-tank storage, friction laws, shaft acceleration, and generator algebra.

\subsection*{Composition verification}
The top-level model checks that all eleven named systems are present and connected through public interfaces.

\subsection*{Structural compilation}
The hierarchical plant is passed through \texttt{mtkcompile}, transforming the connected equations into a reduced numerical system.

\subsection*{Numerical equilibrium test}
The complete model is solved and checked near
\[
f=50\;\mathrm{Hz},
\quad
Q=10\;\mathrm{m^3/s},
\quad
H_s=120\;\mathrm{m},
\quad
y=1.0,
\]
with
\[
P_e\approx8.65242\;\mathrm{MW}.
\]

\section{What this commit enables}
The package now has one equation-based nonlinear model spanning hydraulics, mechanics, electrical dynamics, measurement, and primary control.

The immediate study sequence is:
\begin{enumerate}
\item apply a gate, reference, or load disturbance after a steady interval;
\item inspect headrace and penstock flow transients;
\item inspect surge-tank head oscillation;
\item inspect turbine mechanical and generator electrical power;
\item inspect shaft frequency and governor response;
\item add actuator limits and transient droop;
\item introduce plant-specific parameters and turbine characteristic maps;
\item add analysis points for linearization and frequency-domain analysis;
\item extend the control hierarchy toward AGC and FCR studies.
\end{enumerate}

\section{Current model boundaries}
This is a reduced baseline rather than the final high-fidelity plant model. Planned extensions include hydraulic losses at the operating point, turbine hill-chart interpolation, gate saturation and rate limits, elastic penstock/water hammer, detailed synchronous-machine voltage/reactive-power dynamics, and plant-specific governor structure.

The value of this commit is that these refinements can now be added one subsystem at a time without changing the complete-system workflow.

\section{Conclusion}
This commit establishes the first complete ModelingToolkit hydropower SMIB in OpenHPLjl. It converts the package from a collection of reusable physical components into a connected, initialized, compilable, and simulatable nonlinear plant model with closed-loop droop control. The architecture is now suitable as the base for transient studies, control development, system identification, FCR workflows, and later plant-specific validation.

\end{document}
